\chapter{2006年天津卷（文）}
\section{单选题}
\begin{problem}
已知集合 \(A = \{x \mid -3 \leqslant x \leqslant 1\}\)，\(B = \{x \mid |x| \leqslant 2\}\)，则 \(A \cap B =\)（\quad）

\choices
    {\(\{x\mid -2\leqslant x\leqslant 1\}\)}
    {\(\{x\mid 0\leqslant x\leqslant 1\}\)}
    {\(\{x\mid -3\leqslant x\leqslant 2\}\)}
    {\(\{x\mid 1\leqslant x\leqslant 2\}\)}
\end{problem}

\begin{answer}
A
\end{answer}

\begin{solution}
由 \(\lvert x\rvert\leqslant2\) 得 \(-2\leqslant x\leqslant2\)，所以 \(B=\{x\mid -2\leqslant x\leqslant2\}\). 与 \(A=\{x\mid -3\leqslant x\leqslant1\}\) 取交集，得 \(A\cap B=\{x\mid -2\leqslant x\leqslant1\}\). 因此选 A.
\end{solution}

\begin{problem}
设 \(\{a_n\}\) 是等差数列，\(a_1 + a_3 + a_5 = 9\)，\(a_6 = 9\)，则这个数列的前 \(6\) 项和等于（\quad）

\choices
    {\(12\)}
    {\(24\)}
    {\(36\)}
    {\(48\)}
\end{problem}

\begin{answer}
B
\end{answer}

\begin{solution}
设等差数列的公差为 \(d\). 由等差数列中间项的性质，\(a_1+a_3+a_5=3a_3=9\)，所以 \(a_3=3\). 又 \(a_6=a_3+3d=9\)，得 \(d=2\)，从而 \(a_1=a_3-2d=-1\). 因此前 \(6\) 项和为 \(S_6=\dfrac{6(a_1+a_6)}{2}=3\times8=24\)，故选 B.
\end{solution}

\begin{problem}
设变量 \(x\)，\(y\) 满足约束条件 \(\begin{cases}
y \leqslant x,\\
x + y \geqslant 2,\\
y \geqslant 3x - 6.
\end{cases}\)，则目标函数 \(z = 2x + y\) 的最小值为（\quad）

\choices
    {\(2\)}
    {\(3\)}
    {\(4\)}
    {\(9\)}
\end{problem}

\begin{answer}
B
\end{answer}

\begin{solution}
约束区域的边界直线为 \(y=x\)，\(y=2-x\) 和 \(y=3x-6\). 三条边界直线两两相交得到的点分别为 \((1,1)\)，\((2,0)\) 和 \((3,3)\)，这三点均满足全部约束条件；又由 \(y\leqslant x\) 与 \(y\geqslant2-x\) 得 \(x\geqslant1\)，由 \(y\leqslant x\) 与 \(y\geqslant3x-6\) 得 \(x\leqslant3\)，所以可行域就是以这三点为顶点的三角形. 线性目标函数在多边形可行域上的最小值在顶点处取得. 计算得 \(z(1,1)=3\)，\(z(2,0)=4\)，\(z(3,3)=9\)，故最小值为 \(3\)，选 B.
\end{solution}

\begin{problem}
设 \(P = \log_2 3\)，\(Q = \log_3 2\)，\(R = \log_2(\log_3 2)\)，则（\quad）

\choices
    {\(R < Q < P\)}
    {\(P < R < Q\)}
    {\(Q < R < P\)}
    {\(R < P < Q\)}
\end{problem}

\begin{answer}
A
\end{answer}

\begin{solution}
因为 \(3>2>1\)，所以 \(P=\log_2 3>1\). 又因为 \(1<2<3\)，所以 \(0<Q=\log_3 2<1\)，从而 \(R=\log_2 Q<0\). 因此 \(R<Q<P\)，故选 A.
\end{solution}

\begin{problem}
设 \(\alpha\)，\(\beta \in \left(-\frac{\pi}{2}, \frac{\pi}{2}\right)\)，那么“\(\alpha < \beta\)”是“\(\tan \alpha < \tan \beta\)”的（\quad）

\choices
    {充分而不必要条件}
    {必要而不充分条件}
    {充分必要条件}
    {既不充分也不必要条件}
\end{problem}

\begin{answer}
C
\end{answer}

\begin{solution}
在区间 \(\left(-\dfrac{\pi}{2},\dfrac{\pi}{2}\right)\) 内，函数 \(y=\tan x\) 的导数为 \(y'=\sec^2x>0\)，所以它在该区间上严格递增. 因而 \(\alpha<\beta\) 当且仅当 \(\tan\alpha<\tan\beta\)，故这是充分必要条件，选 C.
\end{solution}

\begin{problem}
函数 \(y = \sqrt{x^2 + 1} + 1\)（\(x < 0\)）的反函数是（\quad）

\choices
    {\(y = \sqrt{x^2 - 2x}\)（\(x < 0\)）}
    {\(y = -\sqrt{x^2 - 2x}\)（\(x < 0\)）}
    {\(y = \sqrt{x^2 - 2x}\)（\(x > 2\)）}
    {\(y = -\sqrt{x^2 - 2x}\)（\(x > 2\)）}
\end{problem}

\begin{answer}
D
\end{answer}

\begin{solution}
原函数的定义域为 \(x<0\)，且 \(y=\sqrt{x^2+1}+1>2\). 由 \(y-1=\sqrt{x^2+1}\) 得 \(x^2=(y-1)^2-1=y^2-2y\). 因为原定义域要求 \(x<0\)，所以 \(x=-\sqrt{y^2-2y}\)，其中 \(y>2\). 交换 \(x\) 与 \(y\)，反函数为 \(y=-\sqrt{x^2-2x}\)，定义域为 \(x>2\)，故选 D.
\end{solution}

\begin{problem}
若 \(l\) 为一条直线，\(\alpha\)，\(\beta\)，\(\gamma\) 为三个互不重合的平面，给出下面三个命题：

\begin{circlelist}
\item \(\alpha \perp \gamma\)，\(\beta \perp \gamma \Rightarrow \alpha \perp \beta\)；
\item \(\alpha \perp \gamma\)，\(\beta \parallel \gamma \Rightarrow \alpha \perp \beta\)；
\item \(l \perp \alpha\)，\(l \perp \beta \Rightarrow \alpha \perp \beta\).
\end{circlelist}

其中结论是（\quad）

\choices
    {\(0\)个}
    {\(1\)个}
    {\(2\)个}
    {\(3\)个}
\end{problem}

\begin{answer}
B
\end{answer}

\begin{solution}
命题\(\circled{1}\)不成立：取 \(\gamma\) 为水平面，取 \(\alpha\)，\(\beta\) 为两个互不垂直的竖直平面，则二者都垂直于 \(\gamma\)，但 \(\alpha\) 与 \(\beta\) 不垂直. 命题\(\circled{2}\)成立，因为一个平面若垂直于某平面，也垂直于与该平面平行的任一平面，所以由 \(\alpha\perp\gamma\) 和 \(\beta\parallel\gamma\) 得 \(\alpha\perp\beta\). 命题\(\circled{3}\)不成立：若一条直线同时垂直于两个平面，则这两个平面的法向量都与该直线平行，因而两个平面平行；在它们互不重合时不可能互相垂直. 所以只有命题\(\circled{2}\)正确，结论有 \(1\) 个，选 B.
\end{solution}

\begin{problem}
椭圆的中心为点 \(E(-1,0)\)，它的一个焦点为 \(F(-3,0)\)，相应于焦点 \(F\) 的准线方程为 \(x = -\frac{7}{2}\)，则这个椭圆的方程是（\quad）

\choices
    {\(\frac{2(x - 1)^2}{21} + \frac{2y^2}{3} = 1\)}
    {\(\frac{2(x + 1)^2}{21} + \frac{2y^2}{3} = 1\)}
    {\(\frac{(x - 1)^2}{5} + y^2 = 1\)}
    {\(\frac{(x + 1)^2}{5} + y^2 = 1\)}
\end{problem}

\begin{answer}
D
\end{answer}

\begin{solution}
椭圆中心为 \((-1,0)\)，焦点为 \((-3,0)\)，所以焦距参数 \(c=2\)，长轴在 \(x\) 轴上. 设长半轴为 \(a\)，则左焦点对应的准线为 \(x=-1-\dfrac{a^2}{c}=-1-\dfrac{a^2}{2}\). 由题设 \(-1-\dfrac{a^2}{2}=-\dfrac{7}{2}\)，得 \(a^2=5\). 再由 \(b^2=a^2-c^2=5-4=1\)，所以椭圆方程为 \(\dfrac{(x+1)^2}{5}+y^2=1\)，故选 D.
\end{solution}

\begin{problem}
已知函数 \(f(x) = a\sin x - b\cos x\)（\(a\)，\(b\) 为常数，\(a \neq 0\)，\(x \in \R\)）的图象关于直线 \(x = \frac{\pi}{4}\) 对称，则函数 \(y = f\left(\frac{3\pi}{4} - x\right)\) 是（\quad）

\choices
    {偶函数且它的图象关于点 \((\pi, 0)\) 对称}
    {偶函数且它的图象关于点 \(\left(\frac{3\pi}{2}, 0\right)\) 对称}
    {奇函数且它的图象关于点 \(\left(\frac{3\pi}{2}, 0\right)\) 对称}
    {奇函数且它的图象关于点 \((\pi, 0)\) 对称}
\end{problem}

\begin{answer}
D
\end{answer}

\begin{solution}
图象关于直线 \(x=\dfrac{\pi}{4}\) 对称，等价于对任意 \(x\)，有 \(f\left(\dfrac{\pi}{2}-x\right)=f(x)\). 代入函数式得 \(a\cos x-b\sin x=a\sin x-b\cos x\). 令 \(x=0\)，得到 \(a=-b\)，所以 \(b=-a\)，从而 \(f(x)=a(\sin x+\cos x)\). 因此

\[
f\left(\frac{3\pi}{4}-x\right)=a\left(\sin\left(\frac{3\pi}{4}-x\right)+\cos\left(\frac{3\pi}{4}-x\right)\right)=\sqrt{2}a\sin x
\]

记此函数为 \(g(x)\)，则 \(g(-x)=-g(x)\)，所以它是奇函数；又 \(g(2\pi-x)=-g(x)\)，故其图象关于点 \((\pi,0)\) 对称. 选 D.
\end{solution}

\begin{problem}
如果函数 \(f(x) = a^x(a^x - 3a^2 - 1)\)（\(a > 0\) 且 \(a \neq 1\)）在区间 \([0,+\infty)\) 上是增函数，那么实数 \(a\) 的取值范围是（\quad）

\choices
    {\(\left(0, \frac{2}{3}\right]\)}
    {\(\left[\frac{\sqrt{3}}{3}, 1\right)\)}
    {\((1, \sqrt{3}]\)}
    {\(\left[\frac{3}{2}, +\infty\right)\)}
\end{problem}

\begin{answer}
B
\end{answer}

\begin{solution}
求导得 \(f'(x)=a^x\ln a\left(2a^x-3a^2-1\right)\). 当 \(0<a<1\) 时，\(\ln a<0\)，且 \(a^x\) 在 \([0,+\infty)\) 上的最大值为 \(1\). 因而要使 \(f'(x)\geqslant0\) 对所有 \(x\geqslant0\) 成立，需且只需 \(2-(3a^2+1)\leqslant0\)，即 \(a\geqslant\dfrac{\sqrt{3}}{3}\). 结合 \(0<a<1\)，得到 \(a\in\left[\dfrac{\sqrt{3}}{3},1\right)\).

当 \(a>1\) 时，\(\ln a>0\)，且 \(a^x\geqslant1\)，要使导数恒非负需 \(2-(3a^2+1)\geqslant0\)，即 \(a\leqslant\dfrac{\sqrt{3}}{3}\)，这与 \(a>1\) 矛盾. 所以实数 \(a\) 的取值范围为 \(\left[\dfrac{\sqrt{3}}{3},1\right)\)，选 B.
\end{solution}

\section{填空题}
\begin{problem}
\(\left(x + \frac{1}{\sqrt{x}}\right)^7\) 的二项展开式中 \(x\) 的系数是\fillinblank{}.（用数字作答）
\end{problem}

\begin{answer}
\(35\)
\end{answer}

\begin{solution}
展开式的第 \(k+1\) 项为 \(\mathrm C_7^k x^{7-k}\left(x^{-1/2}\right)^k=\mathrm C_7^k x^{7-\frac{3k}{2}}\)，其中 \(k=0\)，\(1\)，\(\dots\)，\(7\). 要得到 \(x\) 的一次项，需满足 \(7-\dfrac{3k}{2}=1\)，解得 \(k=4\). 因此所求系数为 \(\mathrm C_7^4=35\).
\end{solution}

\begin{problem}
设向量 \(\bs{a}\) 与 \(\bs{b}\) 的夹角为 \(\theta\)，且 \(\bs{a} = (3, 3)\)，\(2\bs{b} - \bs{a} = (-1, 1)\)，则 \(\cos \theta =\fillinblank{}\).
\end{problem}

\begin{answer}
\(\dfrac{3\sqrt{10}}{10}\)
\end{answer}

\begin{solution}
由 \(2\bs{b}-\bs{a}=(-1,1)\) 和 \(\bs{a}=(3,3)\)，得 \(2\bs{b}=(2,4)\)，所以 \(\bs{b}=(1,2)\). 于是 \(\bs{a}\cdot\bs{b}=3\times1+3\times2=9\)，\(\lvert\bs{a}\rvert=\sqrt{3^2+3^2}=3\sqrt{2}\)，\(\lvert\bs{b}\rvert=\sqrt{1^2+2^2}=\sqrt{5}\). 根据数量积公式，\(\cos\theta=\dfrac{\bs{a}\cdot\bs{b}}{\lvert\bs{a}\rvert\lvert\bs{b}\rvert}=\dfrac{9}{3\sqrt{10}}=\dfrac{3\sqrt{10}}{10}\).
\end{solution}

\begin{problem}
如图，在正三棱柱 \(ABC - A_{1}B_{1}C_{1}\) 中，\(AB = 1\). 若二面角 \(C - AB - C_{1}\) 的大小为 \(60^{\circ}\)，则点 \(C\) 到平面 \(ABC_{1}\) 的距离为\fillinblank{}.

\bitmapfigure[max width=0.32\linewidth]{img/2006/tianjin_liberal/q13_fig1.png}
\end{problem}

\begin{answer}
\(\dfrac{3}{4}\)
\end{answer}

\begin{solution}
建立直角坐标系，令 \(A=(0,0,0)\)，\(B=(1,0,0)\)，\(C=\left(\dfrac12,\dfrac{\sqrt3}{2},0\right)\)，并设正三棱柱的高为 \(h\)，则 \(C_1=\left(\dfrac12,\dfrac{\sqrt3}{2},h\right)\). 平面 \(ABC\) 的一个法向量为 \(\bs n_1=(0,0,1)\)，平面 \(ABC_1\) 的一个法向量可取 \(\bs n_2=\left(0,-h,\dfrac{\sqrt3}{2}\right)\). 二面角 \(C-AB-C_1\) 是这两个平面的锐角，故

\[
\cos60^{\circ}=\frac{\lvert\bs n_1\cdot\bs n_2\rvert}{\lvert\bs n_1\rvert\lvert\bs n_2\rvert}=\frac{\sqrt3/2}{\sqrt{h^2+3/4}}
\]

由此得 \(h^2=\dfrac94\)，所以 \(h=\dfrac32\). 平面 \(ABC_1\) 的方程为 \(-hy+\dfrac{\sqrt3}{2}z=0\)，于是点 \(C\) 到该平面的距离为

\[
\frac{\left|-h\dfrac{\sqrt3}{2}\right|}{\sqrt{h^2+3/4}}=\frac{(3/2)(\sqrt3/2)}{\sqrt3}=\frac34
\]

故所求距离为 \(\dfrac34\).
\end{solution}

\begin{problem}
若半径为 \(1\) 的圆分别与 \(y\) 轴的正半轴和射线 \(y = \frac{\sqrt{3}}{3}x\)（\(x \geqslant 0\)）相切，则这个圆的方程为\fillinblank{}.
\end{problem}

\begin{answer}
\((x-1)^2+(y-\sqrt{3})^2=1\)
\end{answer}

\begin{solution}
设圆心为 \(T(u,v)\). 圆与 \(y\) 轴的正半轴相切，故 \(\lvert u\rvert=1\) 且 \(v>0\). 射线所在直线为 \(x-\sqrt3y=0\)，圆心到该直线的距离为半径 \(1\)，所以 \(\lvert u-\sqrt3v\rvert=2\).

若 \(u=1\)，结合 \(v>0\) 可得 \(v=\sqrt3\). 若 \(u=-1\)，结合 \(v>0\) 得 \(v=\dfrac1{\sqrt3}\)，但此时圆心到射线方向单位向量 \(\left(\dfrac{\sqrt3}{2},\dfrac12\right)\) 的有向投影为 \((-1,1/\sqrt3)\cdot(\sqrt3/2,1/2)=-\sqrt3/3<0\)，垂足落在射线的反向延长线上，不符合相切于该射线的条件. 因而圆心只能是 \((1,\sqrt3)\)，圆的方程为 \((x-1)^2+(y-\sqrt3)^2=1\).
\end{solution}

\begin{problem}
某公司一年购买某种货物 \(400\text{ 吨}\)，每次都购买 \(x\text{ 吨}\)，运费为 \(4\text{ 万元/次}\)，一年的总存储费用为 \(4x\text{ 万元}\)，要使一年的总运费与总存储费用之和最小，则 \(x=\)\fillinblank{}\(\text{吨}\).
\end{problem}

\begin{answer}
\(20\)
\end{answer}

\begin{solution}
取货量须满足 \(x>0\)，一年需要购买 \(\dfrac{400}{x}\) 次，所以一年的总运费与总存储费用之和为 \(T(x)=4\cdot\dfrac{400}{x}+4x=\dfrac{1600}{x}+4x\). 由基本不等式，\(T(x)\geqslant2\sqrt{\dfrac{1600}{x}\cdot4x}=160\)，当且仅当 \(\dfrac{1600}{x}=4x\) 时取等号. 因为 \(x>0\)，解得 \(x=20\)，故每次购买 \(20\text{ 吨}\) 时总费用最小.
\end{solution}

\begin{problem}
用数字 \(0\)，\(1\)，\(2\)，\(3\)，\(4\) 组成没有重复数字的五位数，则其中数字 \(1\)，\(2\) 相邻的偶数有\fillinblank{}\(\text{个}\).（用数字作答）
\end{problem}

\begin{answer}
\(24\)
\end{answer}

\begin{solution}
五位数必须使用这五个数字各一次，且末位只能是 \(0\)，\(2\) 或 \(4\). 将相邻的 \(1\)，\(2\) 看作一个整体时，该整体有 \(12\) 和 \(21\) 两种排列.

末位为 \(0\) 时，前四位由 \(12\)（或 \(21\)）、\(3\)、\(4\) 排列，共有 \(2\times3!=12\) 个. 末位为 \(2\) 时，为保证 \(1\)，\(2\) 相邻，第四位必须为 \(1\)，前三位排列 \(0\)，\(3\)，\(4\)，且首位不能为 \(0\)，共有 \(3!-2!=4\) 个. 末位为 \(4\) 时，前四位的排列中相邻的 \(1\)，\(2\) 共有 \(2\times3!=12\) 个，其中首位为 \(0\) 的有 \(2\times2!=4\) 个，所以有 \(8\) 个.

因此满足条件的五位数共有 \(12+4+8=24\) 个.
\end{solution}

\section{解答题}
\begin{problem}
已知 \(\tan \alpha + \cot \alpha = \frac{5}{2}\)，\(\alpha \in \left(\frac{\pi}{4}, \frac{\pi}{2}\right)\)，求 \(\cos 2\alpha\) 和 \(\sin\left(2\alpha + \frac{\pi}{4}\right)\) 的值.
\end{problem}

\begin{solution}
令 \(t=\tan\alpha\). 因为 \(\alpha\in\left(\frac{\pi}{4},\frac{\pi}{2}\right)\)，所以 \(t>1\)，且 \(\cot\alpha=\frac{1}{t}\). 由已知得

\[
t+\frac{1}{t}=\frac{5}{2}\Longrightarrow 2t^{2}-5t+2=0\Longrightarrow (2t-1)(t-2)=0
\]

结合 \(t>1\)，可得 \(t=2\). 因此 \(\cos 2\alpha=\frac{1-t^{2}}{1+t^{2}}=-\frac{3}{5}\)，且
\(\sin 2\alpha=\frac{2t}{1+t^{2}}=\frac{4}{5}\).

所以

\[
\sin\left(2\alpha+\frac{\pi}{4}\right)=\sin 2\alpha\cos\frac{\pi}{4}+\cos 2\alpha\sin\frac{\pi}{4}=\frac{4}{5}\cdot\frac{\sqrt{2}}{2}-\frac{3}{5}\cdot\frac{\sqrt{2}}{2}=\frac{\sqrt{2}}{10}
\]
\end{solution}

\begin{problem}
甲，乙两台机床相互没有影响地生产某种产品，甲机床产品的正品率是 \(0.9\)，乙机床产品的正品率是 \(0.95\).

\begin{enumerate}
\item 从甲机床生产的产品中任取 \(3\) 件，求其中恰有 \(2\) 件正品的概率（用数字作答）；
\item 从甲，乙两台机床生产的产品中各任取 \(1\) 件，求其中至少有 \(1\) 件正品的概率（用数字作答）.
\end{enumerate}
\end{problem}

\begin{solution}
\begin{enumerate}
\item 甲机床生产的产品中，正品率为 \(0.9\)，次品率为 \(1-0.9=0.1\). 由于各件产品相互独立，恰有 \(2\) 件正品的概率为
\(\mathrm C_{3}^{2}(0.9)^{2}(0.1)=3\times 0.81\times 0.1=0.243\).

\item 两台机床各取出的产品相互独立，至少有 \(1\) 件正品的对立事件是两件都是次品. 因此所求概率为
\(1-(1-0.9)(1-0.95)=1-0.1\times 0.05=0.995\).
\end{enumerate}
\end{solution}

\begin{problem}
如图，在五面体 \(ABCDEF\) 中，点 \(O\) 是矩形 \(ABCD\) 的对角线的交点，面 \(CDE\) 是等边三角形，棱 \(EF \parallel BC\) 且 \(EF = \frac{1}{2}BC\).

\begin{enumerate}
\item 证明：\(FO \parallel\) 平面 \(CDE\)；
\item 设 \(BC = \sqrt{3}CD\)，证明：\(EO \perp\) 平面 \(CDF\).
\end{enumerate}

\FigureLayoutDeclare{img/2006/tianjin_liberal/q19_fig1.png}{7.0cm}{45bp 59bp 43bp 61bp}
\bitmapfigure[max width=0.44\linewidth]{img/2006/tianjin_liberal/q19_fig1.png}
\end{problem}

\begin{solution}
取 \(M\) 为 \(CD\) 的中点，连接 \(OM\)、\(EM\)、\(FM\). 因为 \(ABCD\) 是矩形，且 \(O\) 是其对角线的交点，所以 \(OM\parallel BC\)，\(OM=\dfrac{1}{2}BC\). 又由已知 \(EF\parallel BC\)，\(EF=\dfrac{1}{2}BC\)，可得 \(EF\parallel MO\) 且 \(EF=MO\). 因此四边形 \(EFOM\) 是平行四边形.

\begin{enumerate}
\item 由平行四边形的对边平行，得 \(FO\parallel EM\). 直线 \(EM\) 在平面 \(CDE\) 内. 又因为平面 \(ABCD\) 与平面 \(CDE\) 的交线为 \(CD\)，而 \(O\notin CD\)，所以 \(O\notin\) 平面 \(CDE\). 因而 \(FO\) 不在平面 \(CDE\) 内，从而 \(FO\parallel\) 平面 \(CDE\).

\item 在等边三角形 \(CDE\) 中，\(EM\perp CD\)，且 \(EM=\dfrac{\sqrt{3}}{2}CD\). 由 \(BC=\sqrt{3}CD\) 及 \(EF=\dfrac{1}{2}BC\)，得

\[
EF=\frac{\sqrt{3}}{2}CD=EM
\]

所以平行四边形 \(EFOM\) 是菱形，设其两条对角线的交点为 \(P\)，则 \(P\in EO\cap FM\)，且 \(EO\perp FM\).

矩形的邻边垂直，且 \(OM\parallel BC\)，所以 \(CD\perp OM\). 又 \(CD\perp EM\)，而 \(OM\) 与 \(EM\) 是平面 \(EOM\) 内相交的两条直线，故 \(CD\perp\) 平面 \(EOM\). 因而 \(CD\perp EO\). 又 \(FM\perp EO\)，且 \(CD\) 与 \(FM\) 在点 \(M\) 相交，二者同在平面 \(CDF\) 内，所以 \(EO\perp\) 平面 \(CDF\).
\end{enumerate}
\end{solution}

\begin{problem}
已知函数 \(f(x) = 4x^{3} - 3x^{2}\cos \theta + \frac{1}{32}\)，其中 \(x \in \R\)，\(\theta\) 为参数，且 \(0 \leqslant \theta \leqslant \frac{\pi}{2}\).

\begin{enumerate}
\item 当 \(\cos\theta = 0\) 时，判断函数 \(f(x)\) 是否有极值；
\item 要使函数 \(f(x)\) 的极小值大于零，求参数 \(\theta\) 的取值范围；
\item 若对（2）中所求的取值范围内的任意参数 \(\theta\)，函数 \(f(x)\) 在区间 \((2a - 1, a)\) 内都是增函数，求实数 \(a\) 的取值范围.
\end{enumerate}
\end{problem}

\begin{solution}
\begin{enumerate}
\item 当 \(\cos\theta=0\) 时，\(f(x)=4x^{3}+\frac{1}{32}\)，它在 \(\R\) 上为增函数，因此没有极值.

\item 记 \(c=\cos\theta\)，则 \(0\leqslant c\leqslant 1\)，且 \(f'(x)=12x^{2}-6cx=6x(2x-c)\).

当 \(c>0\) 时，\(f'(x)\) 在 \((-\infty,0)\) 和 \(\left(\frac{c}{2},+\infty\right)\) 上为正，在 \(\left(0,\frac{c}{2}\right)\) 上为负，所以函数在 \(x=\frac{c}{2}\) 处取得极小值. 该极小值为

\(f\left(\frac{c}{2}\right)=4\left(\frac{c}{2}\right)^{3}-3\left(\frac{c}{2}\right)^{2}c+\frac{1}{32}=\frac{1}{32}-\frac{c^{3}}{4}\).

要使极小值大于零，需满足 \(c>0\) 且

\(\frac{1}{32}-\frac{c^{3}}{4}>0\Longleftrightarrow c<\frac{1}{2}\).

结合 \(0\leqslant\theta\leqslant\frac{\pi}{2}\)，可得

\(\frac{\pi}{3}<\theta<\frac{\pi}{2}\).

\item 由（2）知，对任意所取参数，\(c\in\left(0,\frac{1}{2}\right)\). 对固定的 \(c\)，有 \(f'(x)<0\) 当且仅当 \(0<x<\frac{c}{2}\)，而在 \(x<0\) 或 \(x>\frac{c}{2}\) 时 \(f'(x)>0\). 因而区间 \(I=(2a-1,a)\) 对每一个这样的参数都满足函数递增，当且仅当 \(I\) 不与任意一个区间 \(\left(0,\frac{c}{2}\right)\) 相交. 又因为

\[
\bigcup_{0<c<1/2}\left(0,\frac{c}{2}\right)=\left(0,\frac{1}{4}\right)
\]

所以非空区间 \(I\) 不能含有 \(\left(0,\frac{1}{4}\right)\) 中的点. 由于 \(I\) 是连通区间，它必须完全落在 \((-\infty,0]\) 或 \([\frac{1}{4},+\infty)\) 内. 前一种情形给出 \(a\leqslant0\). 后一种情形给出 \(2a-1\geqslant\frac{1}{4}\)，即 \(a\geqslant\frac{5}{8}\)，同时区间非空要求 \(2a-1<a\)，即 \(a<1\). 反之，这两组条件分别使 \(I\) 落在导数为正的区域内，对每个 \(c\in\left(0,\frac{1}{2}\right)\) 都成立. 因而

\(a\in\left(-\infty,0\right]\cup\left[\frac{5}{8},1\right)\).
\end{enumerate}
\end{solution}

\begin{problem}
已知数列 \(\{x_n\}\) 满足 \(x_1 = x_2 = 1\)，并且 \(\frac{x_{n+1}}{x_n} = \lambda \frac{x_n}{x_{n-1}}\)（\(\lambda\) 为非零参数，\(n = 2\)，\(3\)，\(4\)，\(\dots\)）.

\begin{enumerate}
\item 若 \(x_1\)，\(x_3\)，\(x_5\) 成等比数列，求参数 \(\lambda\) 的值；
\item 设 \(0 < \lambda < 1\)，常数 \(k \in \N^*\) 且 \(k \geqslant 3\)，证明：\(\frac{x_{1+k}}{x_1} + \frac{x_{2+k}}{x_2} + \cdots + \frac{x_{n+k}}{x_n} < \frac{\lambda^k}{1 - \lambda^k}\)（\(n \in \N^*\)）.
\end{enumerate}
\end{problem}

\begin{solution}
由 \(x_1=x_2=1\ne0\) 且 \(\lambda\ne0\)，根据递推关系可知数列各项均不为零. 令

\(r_n=\frac{x_{n+1}}{x_n}\)，\((n\in\N^*)\).

则 \(r_1=1\)，并且对 \(n\geqslant2\)，有

\[
r_n=\frac{x_{n+1}}{x_n}=\lambda\frac{x_n}{x_{n-1}}=\lambda r_{n-1}
\]

所以 \(r_n=\lambda^{n-1}\)，从而
\(x_n=x_1r_1r_2\cdots r_{n-1}=\lambda^{\frac{(n-1)(n-2)}{2}}\).

\begin{enumerate}
\item 由上式得 \(x_1=1\)，\(x_3=\lambda\)，\(x_5=\lambda^{6}\). 三项成等比数列，当且仅当

\[
x_3^{2}=x_1x_5\Longrightarrow \lambda^{2}=\lambda^{6}\Longrightarrow \lambda^{4}=1
\]

因为 \(\lambda\) 为非零实数，所以 \(\lambda=1\) 或 \(\lambda=-1\).

\item 对任意 \(j\in\N^*\)，有

\[
\frac{x_{j+k}}{x_j}=\lambda^{\frac{(j+k-1)(j+k-2)-(j-1)(j-2)}{2}}=\lambda^{kj+\frac{k(k-3)}{2}}
\]

因此

\[
\begin{gathered}
\sum_{j=1}^{n}\frac{x_{j+k}}{x_j}
=\sum_{j=1}^{n}\lambda^{kj+\frac{k(k-3)}{2}}
=\lambda^{\frac{k(k-1)}{2}}\left(1+\lambda^{k}+\lambda^{2k}+\cdots+\lambda^{(n-1)k}\right)\\
=\frac{\lambda^{\frac{k(k-1)}{2}}\left(1-\lambda^{kn}\right)}{1-\lambda^{k}}
\end{gathered}
\]

因为 \(0<\lambda<1\) 且 \(k\geqslant3\)，所以 \(\frac{k(k-1)}{2}\geqslant k\)，从而 \(\lambda^{\frac{k(k-1)}{2}}\leqslant\lambda^{k}\)，又 \(0<1-\lambda^{kn}<1\). 因此

\[
\sum_{j=1}^{n}\frac{x_{j+k}}{x_j}<\frac{\lambda^{k}}{1-\lambda^{k}}
\]
\end{enumerate}
\end{solution}

\begin{problem}
如图，双曲线 \(\frac{x^2}{a^2} - \frac{y^2}{b^2} = 1\)（\(a > 0\)，\(b > 0\)）的离心率为 \(\frac{\sqrt{5}}{2}\)，\(F_1\)，\(F_2\) 分别为左，右焦点，\(M\) 为左准线与渐近线在第二象限内的交点，且 \(\overrightarrow{F_1M} \cdot \overrightarrow{F_2M} = -\frac{1}{4}\).

\begin{enumerate}
\item 求双曲线的方程；
\item 设 \(A(m,0)\) 和 \(B\left(\frac{1}{m},0\right)\)（\(0 < m < 1\)）是 \(x\) 轴上的两点，过点 \(A\) 作斜率不为 \(0\) 的直线 \(l\)，使得 \(l\) 交双曲线于 \(C\)，\(D\) 两点，作直线 \(BC\) 交双曲线于另一点 \(E\). 证明：直线 \(DE\) 垂直于 \(x\) 轴.
\end{enumerate}

\FigureLayoutDeclare{img/2006/tianjin_liberal/q22_fig1.png}{7.7cm}{156bp 17bp 63bp 31bp}
\bitmapfigure[max width=0.42\linewidth]{img/2006/tianjin_liberal/q22_fig1.png}
\end{problem}

\begin{solution}
\begin{enumerate}
\item 设双曲线的半焦距为 \(c\)，则

\(c^{2}=a^{2}+b^{2}\)，\(\frac{c}{a}=\frac{\sqrt{5}}{2}\).

所以

\(c^{2}=\frac{5}{4}a^{2}\)，\(b^{2}=c^{2}-a^{2}=\frac{1}{4}a^{2}\).

左准线方程为 \(x=-\frac{a^{2}}{c}\)，渐近线方程为 \(y=\pm\frac{b}{a}x=\pm\frac{x}{2}\). 因为 \(M\) 在第二象限，所以

\(M\left(-\frac{a^{2}}{c},\frac{a^{2}}{2c}\right)\)，\(F_1(-c,0)\)，\(F_2(c,0)\).

于是

\(\overrightarrow{F_1M}=\left(\frac{c^{2}-a^{2}}{c},\frac{a^{2}}{2c}\right)\)，
\(\overrightarrow{F_2M}=\left(-\frac{a^{2}+c^{2}}{c},\frac{a^{2}}{2c}\right)\).

由 \(c^{2}=\frac{5}{4}a^{2}\) 得

\[
\overrightarrow{F_1M}\cdot\overrightarrow{F_2M}=-\frac{c^{4}-a^{4}}{c^{2}}+\frac{a^{4}}{4c^{2}}=-\frac{5a^{4}}{16c^{2}}=-\frac{a^{2}}{4}
\]

结合题设 \(\overrightarrow{F_1M}\cdot\overrightarrow{F_2M}=-\frac{1}{4}\)，可得 \(a^{2}=1\). 因为 \(a>0\)，所以 \(a=1\)，\(b=\frac{1}{2}\). 双曲线方程为

\[
x^{2}-4y^{2}=1
\]

\item 由（1）知双曲线为 \(x^{2}-4y^{2}=1\). 设直线 \(l\) 的斜率为 \(r\)，令

\(C=(u,r(u-m))\)，\(D=(v,r(v-m))\).

其中 \(u\)，\(v\) 是方程

\[
x^{2}-4r^{2}(x-m)^{2}=1
\]

的两个不同实根. 因为直线与双曲线有两个不同的有限交点，所以 \(4r^{2}\ne1\). 由韦达定理，

\[
u+v=\frac{8mr^{2}}{4r^{2}-1}
\]

设直线 \(BC\) 的斜率为 \(s\)，则

若 \(u=\frac{1}{m}\)，由点 \(C\) 在双曲线上可得

\(4r^{2}=\frac{m^{-2}-1}{(m^{-1}-m)^{2}}=\frac{1}{1-m^{2}}\).

于是

\(u+v=\frac{8mr^{2}}{4r^{2}-1}=\frac{2}{m}=2u\)，即 \(v=u\)，这与 \(C\)，\(D\) 是两个不同交点矛盾. 所以 \(u\ne\frac{1}{m}\). 因而
\(s=\frac{r(u-m)}{u-1/m}\).

点 \(E\) 的横坐标记为 \(w\). 由点 \(C\) 在双曲线上，

\(4r^{2}=\frac{u^{2}-1}{(u-m)^{2}}\).

从而

\[
4s^{2}=\frac{4r^{2}(u-m)^{2}}{(u-1/m)^{2}}=\frac{m^{2}(u^{2}-1)}{(mu-1)^{2}}
\]

且

\[
4s^{2}-1=\frac{2mu-m^{2}-1}{(mu-1)^{2}}
\]

因为直线 \(BC\) 与双曲线有两个不同的有限交点 \(C\)，\(E\)，所以 \(4s^{2}\ne1\). 直线 \(BC\) 与双曲线的两个交点横坐标为 \(u\)，\(w\)，仍由韦达定理得

\(u+w=\frac{8s^{2}/m}{4s^{2}-1}=\frac{2m(u^{2}-1)}{2mu-m^{2}-1}\).

另一方面，由 \(4r^{2}=\frac{u^{2}-1}{(u-m)^{2}}\) 可得
\(4r^{2}-1=\frac{2mu-m^{2}-1}{(u-m)^{2}}\).

所以

\[
u+v=\frac{8mr^{2}}{4r^{2}-1}=\frac{2m(u^{2}-1)}{2mu-m^{2}-1}=u+w
\]

于是 \(v=w\)，即 \(x_D=x_E\). 若 \(D=E\)，则直线 \(BC\) 经过 \(C\)，\(D\)，从而与直线 \(l=CD\) 重合，导致 \(B\in l\)，这与直线 \(l\) 斜率不为零矛盾. 因此 \(D\ne E\)，直线 \(DE\) 为竖直直线，故 \(DE\perp x\) 轴.
\end{enumerate}
\end{solution}
